% Partie : intuitionistic-logic
% Chapitre : tableaux
% Section : introduction

\documentclass[../../../include/open-logic-section]{subfiles}

\begin{document}

\olfileid[fr]{int}{tab}{int}

\olsection{Introduction}

Les !!^{tableau}s sont des arbres ramifiés vers le bas dont les nœuds
portent des !!{signed formula}s, c'est-à-dire des couples formés d'un
signe de vérité ($\True$ ou $\False$) et d'!!a{sentence} :
\[
\sFmla{\True}{!A} \text{ ou } \sFmla{\False}{!A}.
\]
!!^a{tableau} commence par un certain nombre d'\emph{hypothèses}.
Chaque autre !!{signed formula} résulte de l'application d'une règle
d'inférence. Certaines règles ajoutent une ou plusieurs !!{signed
formula}s à l'extrémité d'une branche ; d'autres ajoutent deux
nouvelles extrémités et créent ainsi deux branches. Les règles de
décomposition produisent des !!{signed formula}s dont la !!{formula}
est moins complexe que celle de la !!{signed formula} à laquelle elles
sont appliquées. Lorsqu'une branche contient à la fois
$\sFmla{\True}{!A}$ et $\sFmla{\False}{!A}$, on dit qu'elle est
\emph{fermée}. Si toutes les branches d'!!a{tableau} sont fermées,
le !!{tableau} entier est fermé. Un !!{tableau} fermé constitue
!!a{derivation} montrant que l'ensemble des !!{signed formula}s qui
ont servi à commencer le !!{tableau} n'est pas satisfaisable. Cela
permet de définir une relation $\Proves$ : $\Gamma \Proves !A$ si et
seulement s'il existe un ensemble fini~$\Gamma_0 = \{!B_1,
\dots, !B_n\} \subseteq \Gamma$ tel qu'un !!{tableau} fermé ait
pour hypothèses
\[
\{\sFmla{\False}{!A}, \sFmla{\True}{!B_1}, \dots, \sFmla{\True}{!B_n}\}.
\]

Pour la logique intuitionniste, il faut étendre la notion de
!!{signed formula} et adapter les règles des connecteurs. Outre un
signe ($\True$ ou $\False$), les !!{formula}s des !!{tableau}s
intuitionnistes portent des \emph{préfixes}~$\sigma$. Ces préfixes
sont des suites non vides d'entiers strictement positifs : $\sigma
\in (\PosInt)^* \setminus \{\emptyseq\}$. Nous les écrivons sans
les chevrons $\tuple{\ }$ et séparons leurs !!{element}s par des
points~$.$ plutôt que par des virgules $,$. Si $\sigma$ est un
préfixe, $\sigma.n$ désigne $\sigma \concat \tuple{n}$ ; par
exemple, si $\sigma = 1.2.1$, alors $\sigma.3$ est $1.2.1.3$.
Ainsi,
\[
\sFmla{\True}{!A \lif (!B \lif !C)}[1.2]
\]
est une \emph{!!{signed formula} préfixée} (ou, plus brièvement, une
\emph{!!{formula} préfixée}).

Intuitivement, le préfixe désigne un monde d'un modèle susceptible de
satisfaire les !!{formula}s d'une branche d'!!a{tableau} ; si
$\sigma$ désigne un monde, alors $\sigma.n$ désigne un monde
accessible depuis celui que désigne~$\sigma$.

Dans les modèles intuitionnistes, la relation d'accessibilité est
réflexive et transitive. Dans le tableau, $\sigma.*$ désigne
le préfixe~$\sigma$ lui-même ou une extension de la forme
$\sigma.n_1.\cdots.n_k$. Par réflexivité et transitivité, chacun
de ces préfixes désigne un monde accessible depuis celui que
désigne~$\sigma$. On appelle ces préfixes, dans le tableau,
les préfixes accessibles depuis~$\sigma$.

\end{document}
